%==============================================================================
%  ANNEXE - Pile technologique detaillee et comparatif justifie
%  Implementations concretes du socle retenu + alternatives evaluees
%  A inserer via \input{annexe_technologies} dans cahier_des_charges.tex
%  (utilise la charte "Miel & noir" : miel, mielfonce, encadre, colonne C,
%   longtable, booktabs -- deja definis dans le preambule principal)
%==============================================================================

%==============================================================================
\chapter{Annexe : Pile technologique et comparatif des alternatives}
\label{chap:technologies}
%==============================================================================
Le socle technique de \textbf{Zümm} vise l'\textbf{alignement sur les pratiques
du marché} : API REST documentée, client découplé, identité fédérée, données
géospatiales et séries temporelles dans un SGBD unique. Cette annexe recense les
technologies retenues couche par couche, puis justifie chaque choix face à ses
alternatives.

%------------------------------------------------------------------------------
\section{Pile technologique retenue}
%------------------------------------------------------------------------------
Le tableau ci-dessous précise, pour chaque couche de l'architecture
multi-niveaux (figure~\ref{fig:couches}), l'implémentation concrète retenue et
sa version de référence. Toutes les briques sont \textbf{libres et gratuites},
conformément à l'exigence d'autodéploiement sans abonnement.

\begin{xltabular}{\textwidth}{>{\bfseries}p{3.1cm} p{4.9cm} X}
\toprule
\textbf{Couche / besoin} & \textbf{Technologie retenue} & \textbf{Rôle dans Zümm} \\
\midrule
\endhead
Plateforme        & Spring Boot 3 sur JDK 17 LTS      & Socle applicatif : conteneur IoC, autoconfiguration, Tomcat embarqué \\
Exécution         & Image conteneur (Docker)          & Artefact unique déployable, sans serveur d'applications à administrer \\
Contrôle          & Spring MVC + Bean Validation      & Aiguillage des requêtes, validation des entrées \\
Métier            & Composants Spring \texttt{@Service} & Règles de gestion : ROI, seuils, quantité de miel \\
Accès données     & Spring Data JPA (CRUD) + jOOQ (SQL typé) & Persistance des entités, requêtes analytiques et géospatiales \\
SGBD              & PostgreSQL + PostGIS + TimescaleDB & Persistance relationnelle, géolocalisation des sites, séries de mesures \\
Présentation      & React 19 + TypeScript, jetons Zümm & IHM responsive, parité web / mobile, identité visuelle de la charte \\
Graphiques        & Chart.js                          & Tableaux de bord production, alertes, synthèse \\
Cartographie      & MapLibre GL JS + tuiles OpenStreetMap & Carte des sites et rayons de butinage \\
i18n              & Fichiers de ressources par locale & Textes FR / EN / AR, support du RTL arabe \\
Configuration     & Profils Spring (technique) + \texttt{ConfigZumm.ini} via \texttt{PropertySource} dédié (métier) & Infrastructure par environnement ; seuils, unités et modules sans redéploiement \\
API tierce        & REST + contrat OpenAPI 3          & \texttt{getZummHoneyActualQuantity(int)} \\
Authentification  & Keycloak (OpenID Connect) + fédération Google & Connexion déléguée, RBAC et jetons JWT \\
Sécurité échanges & TLS 1.3 + certificats X.509       & Chiffrement et authentification inter-composants \\
Hors-ligne        & PWA (Service Worker) + moteur de synchronisation & Saisie terrain sans réseau, résolution de conflits au retour du réseau \\
Ingestion capteurs & MQTT (courtier EMQX)             & Réception des mesures des ruches connectées \\
Observabilité     & OpenTelemetry vers Prometheus / Grafana & Traces, métriques et tableaux de bord d'exploitation \\
Build \& tests    & Maven, JUnit 5, Mockito, Testcontainers & Compilation, tests unitaires et d'intégration sur base réelle \\
\bottomrule
\end{xltabular}

\begin{encadre}
\textbf{Note :} le domaine métier et les identifiants publics sont indépendants
du socle technique. Les noms de classes, de tables et d'opérations --- dont
\texttt{getZummHoneyActualQuantity} --- restent inchangés quelle que soit la
technologie qui les porte.
\end{encadre}

%------------------------------------------------------------------------------
\section{Comparatif justifié des alternatives}
%------------------------------------------------------------------------------
Pour chaque décision ouverte, les options crédibles ont été évaluées. Le tableau
retient l'option, cite les alternatives écartées et donne la justification au
regard des exigences du projet (autodéploiement, gratuité, ergonomie,
géolocalisation, interopérabilité).

\subsection{Socle applicatif}
\begin{xltabular}{\textwidth}{>{\bfseries}p{2.7cm} p{4.2cm} X}
\toprule
\textbf{Option} & \textbf{Alternatives écartées} & \textbf{Justification du choix} \\
\midrule
\endhead
Spring Boot 3 & Jakarta EE / WildFly, Quarkus, Micronaut, Node.js (NestJS) & \textbf{Écosystème le plus répandu} du marché Java : documentation, bibliothèques et vivier de développeurs abondants, donc coût de maintenance maîtrisé. Quarkus démarre plus vite et consomme moins, mais son écosystème est plus étroit. Jakarta EE impose un serveur d'applications à administrer pour un bénéfice nul à cette échelle. \\
\bottomrule
\end{xltabular}

\subsection{Accès aux données}
\begin{xltabular}{\textwidth}{>{\bfseries}p{2.7cm} p{4.2cm} X}
\toprule
\textbf{Option} & \textbf{Alternatives écartées} & \textbf{Justification du choix} \\
\midrule
\endhead
JPA + jOOQ & JDBC manuel, MyBatis, JPA seul & Approche \textbf{mixte} : JPA supprime le code répétitif du CRUD sur la quinzaine d'entités du domaine, tandis que jOOQ exprime en SQL typé et vérifié à la compilation les requêtes analytiques et les appels PostGIS que JPA rend illisibles. Le JDBC manuel donnerait le même résultat pour un coût de développement bien supérieur. \\
\bottomrule
\end{xltabular}

\subsection{Système de gestion de base de données}
\begin{xltabular}{\textwidth}{>{\bfseries}p{2.7cm} p{4.2cm} X}
\toprule
\textbf{Option} & \textbf{Alternatives écartées} & \textbf{Justification du choix} \\
\midrule
\endhead
PostgreSQL & MySQL / MariaDB, H2, Oracle XE & \textbf{Géolocalisation} native (PostGIS) pour les sites et rayons de butinage et \textbf{séries temporelles} (TimescaleDB) pour les mesures de capteurs, dans une \emph{seule} instance : ni base spécialisée supplémentaire ni synchronisation entre deux systèmes. Bon support des contraintes \texttt{CHECK} (règles cadres/hausses). Gratuit et sans verrouillage propriétaire, contrairement à Oracle. \\
\bottomrule
\end{xltabular}

\subsection{Style de l'API tierce}
\begin{xltabular}{\textwidth}{>{\bfseries}p{2.7cm} p{4.2cm} X}
\toprule
\textbf{Option} & \textbf{Alternatives écartées} & \textbf{Justification du choix} \\
\midrule
\endhead
REST + OpenAPI 3 & SOAP (JAX-WS), GraphQL, gRPC & Standard \textbf{de fait} pour l'intégration tierce : contrat lisible par l'humain et par la machine, génération automatique de clients dans la plupart des langages, outillage de test universel. SOAP n'aurait de sens que pour s'intégrer à un existant qui l'impose. GraphQL répond à des besoins de requêtage client que Zümm n'a pas. \\
\bottomrule
\end{xltabular}

\subsection{Interface utilisateur}
\begin{xltabular}{\textwidth}{>{\bfseries}p{2.7cm} p{4.2cm} X}
\toprule
\textbf{Option} & \textbf{Alternatives écartées} & \textbf{Justification du choix} \\
\midrule
\endhead
React + TypeScript & JSP / Thymeleaf, Vue, Angular, Svelte & \textbf{Découplage} complet du client et du serveur, indispensable au fonctionnement hors-ligne : la même application sert de PWA sur le terrain. TypeScript sécurise les contrats de données générés depuis OpenAPI. Le rendu serveur (JSP, Thymeleaf) interdirait le mode déconnecté. \\
\bottomrule
\end{xltabular}

\subsection{Authentification et contrôle d'accès}
\begin{xltabular}{\textwidth}{>{\bfseries}p{2.7cm} p{4.2cm} X}
\toprule
\textbf{Option} & \textbf{Alternatives écartées} & \textbf{Justification du choix} \\
\midrule
\endhead
Keycloak (OIDC) & OAuth Google câblé à la main, Auth0, Spring Security seul & Fournit \textbf{nativement} les trois besoins du projet : fédération Google, comptes locaux de repli et matrice de rôles. Recoder ce triplet représenterait plusieurs semaines pour un résultat moins sûr. Auth0 est équivalent mais payant au-delà d'un quota, ce qui contredit l'autodéploiement. \\
\bottomrule
\end{xltabular}

\subsection{Externalisation de la configuration}
\begin{xltabular}{\textwidth}{>{\bfseries}p{2.7cm} p{4.2cm} X}
\toprule
\textbf{Option} & \textbf{Alternatives écartées} & \textbf{Justification du choix} \\
\midrule
\endhead
Approche mixte & Profils Spring seuls, \texttt{ConfigZumm.ini} seul & Les deux configurations n'ont ni le même public ni le même cycle de vie. L'\textbf{infrastructure} (sources de données, points d'accès, secrets) change par environnement et relève de l'exploitation technique : profils Spring et variables d'environnement. Les \textbf{seuils métier} (production, alertes, unités) sont réglés par l'apiculteur lui-même, en cours de saison : ils restent dans \texttt{ConfigZumm.ini}, fichier texte lisible, monté en volume et relu à chaud par un \texttt{PropertySource} dédié. Tout confondre dans \texttt{application.yml} obligerait l'utilisateur métier à passer par un technicien. \\
\bottomrule
\end{xltabular}

\subsection{Cartographie et rayons de butinage}
\begin{xltabular}{\textwidth}{>{\bfseries}p{2.7cm} p{4.2cm} X}
\toprule
\textbf{Option} & \textbf{Alternatives écartées} & \textbf{Justification du choix} \\
\midrule
\endhead
MapLibre GL + OSM & Leaflet, Google Maps JS API, Mapbox & \textbf{Libre et sans clé payante} ni quota facturé, cohérent avec l'exigence d'autodéploiement. Tuiles vectorielles : rendu net à tout niveau de zoom et style dérivable des couleurs de la charte, ce que le raster de Leaflet ne permet pas. \\
\bottomrule
\end{xltabular}

\subsection{Bibliothèque de graphiques}
\begin{xltabular}{\textwidth}{>{\bfseries}p{2.7cm} p{4.2cm} X}
\toprule
\textbf{Option} & \textbf{Alternatives écartées} & \textbf{Justification du choix} \\
\midrule
\endhead
Chart.js & D3.js, ApexCharts, Highcharts & Courbe d'apprentissage faible et rendu direct des histogrammes et courbes des tableaux de bord (production, ROI). D3 est plus puissant mais surdimensionné. Highcharts n'est pas libre pour un usage commercial. \\
\bottomrule
\end{xltabular}

\subsection{Accès hors-ligne et mobile}
\begin{xltabular}{\textwidth}{>{\bfseries}p{2.7cm} p{4.2cm} X}
\toprule
\textbf{Option} & \textbf{Alternatives écartées} & \textbf{Justification du choix} \\
\midrule
\endhead
PWA + moteur de synchronisation & Application native (Android/iOS), Flutter, React~Native, IndexedDB nu & \textbf{Parité web / mobile} avec une seule base de code, saisie terrain sans réseau puis synchronisation différée. Un moteur de synchronisation dédié apporte une \textbf{politique de résolution de conflits} éprouvée --- deux agents éditant la même visite hors réseau est le cas nominal du métier, non un cas limite. \\
\bottomrule
\end{xltabular}

\subsection{Tests d'intégration}
\begin{xltabular}{\textwidth}{>{\bfseries}p{2.7cm} p{4.2cm} X}
\toprule
\textbf{Option} & \textbf{Alternatives écartées} & \textbf{Justification du choix} \\
\midrule
\endhead
Testcontainers & Arquillian, base H2 en mémoire & Les tests s'exécutent sur un \textbf{PostgreSQL réel} muni de PostGIS et TimescaleDB, donc sur les mêmes extensions qu'en production --- ce qu'une base H2 en mémoire ne sait pas simuler. Arquillian suppose un serveur d'applications, que l'architecture retenue n'utilise plus. \\
\bottomrule
\end{xltabular}

\begin{encadre}
\textbf{Cohérence avec les contraintes du projet :} toutes les briques retenues
sont libres, documentées et \emph{justifiables} « pourquoi / comment / où »
(critère \emph{Technologies} de la grille d'évaluation, chapitre des délais).
Aucune ne repose sur un générateur de code : elles sont mises en \oe uvre à la
main par l'équipe, conformément à la contrainte de l'épreuve.
\end{encadre}
